\chapter{Smooth morphisms: extension properties}
\label{III}

\section{Formally smooth homomorphisms}
\label{III.1}

In~\SourceRef{II}{II}, we limited ourselves to the consideration of homomorphisms of finite type, and consequently, in local homomorphisms~$A\to B$ of local rings, to the case where~$B$ is isomorphic to a localization of an $A$-algebra of finite type. This case is insufficient for various applications, notably in formal geometry or analytic geometry. For example, the formal power-series ring~$B=A[[t_1,\dots,t_n]]$ has, from the point of view of formal geometry, the properties of a smooth algebra over~$A$. In analytic geometry, the same is true of the local ring of a point~$(y,z)$ of a product $Y\times\CC^n$ considered as an algebra over the local ring of~$y$; moreover, the completion of this algebra is isomorphic to the algebra of formal power series in~$n$ indeterminates over the completion of the base ring~$\cal{O}_x$. This leads to the following definition.

\begin{definition}
\label{III.1.1}
Let~$u\colon A\to B$ be a local homomorphism of local rings (noetherian, as recalled). Suppose~$\kres(B)$ is finite over~$\kres(A)$. We say that~$u$ is a \emph{formally smooth homomorphism}, or that the algebra~$B$ is \emph{formally smooth over}~$A$, if there exists a local finite~$\overline{A}$-algebra~$A'$, free over~$\overline{A}$, such that the local components of the semi-local ring~$\overline{B}\otimes_{\overline{A}} A'=B'$ are $A'$-isomorphic to algebras of formal power series over~$A'$\footnote{For a more general and more conceptual definition, motivated by~\ref{III.2.1} below, \cf EGA~$\mathrm{0}_{\mathrm{IV}}$~19.3.1.}.
\end{definition}

(We denote by~$\overline{A}$, $\overline{B}$ the completed rings of~$A$, $B$.) Since~$B'$ is finite and free over~$\overline{B}$, it is indeed a semi-local ring, a direct sum of complete local rings, each of which is still a free module over~$\overline{B}$, hence has the same dimension as~$\overline{B}$, and therefore as~$B$. It follows that the number of variables~$t_i$ in the formal power-series rings considered in~\ref{III.1.1} is equal to $\dim \overline{B}-\dim \overline{A}=\dim B-\dim A$, and in particular is independent of the local component chosen. One sees at once that it is also the dimension of the ring $B\otimes k=B/\goth{m} B$, where~$k=A/\goth{m}$ is the residue field of~$A$; we shall call it the \emph{relative dimension of}~$B$ \emph{with respect to}~$A$.

\begin{remarks}
\label{III.1.2}
It is evident that Definition~\ref{III.1.1} depends only on the homomorphism on the completions~$\overline{A}\to\overline{B}$ induced by~$A\to B$, which justifies the terminology to some extent. We repent here of Definitions~I~\SourceRef{I.3.2}{3.2}~b) and I~\SourceRef{I.4.1}{4.1}~b), which risk being misleading, and prefer to say ``formally unramified'' respectively ``formally étale'' in the cases considered in these definitions, reserving the terminology ``unramified'' and ``étale'' for the case where~$B$ is a localization of an $A$-algebra of finite type\footnote{Or better, ``essentially unramified'' respectively ``essentially étale'', compare EGA IV 18.6.1.}. The reader will immediately verify that ``formally étale'' is equivalent to ``formally smooth and quasi-finite''. Finally, let us point out that there is a reasonable definition of ``formally smooth'' without any prior hypothesis on the residual extension
$\kres(B)/\kres(A)$ (supposed finite here), encompassing among others the local homomorphisms~$A\to B$ such that~$B$ is \emph{flat} over~$A$ and $B/\goth{m}B$ is a \emph{separable extension} of~$A/\goth{m}=k$ (not necessarily of finite type); for example, a Cohen $p$-ring is formally smooth over the ring of $p$-adic integers. It is the lifting property for homomorphisms (compare~\ref{III.2.1}) that should become the definition in this general case. For the applications we have in view, the case treated in Definition~\ref{III.1.1} will suffice; in what follows, by ``formally smooth'' we shall understand ``with finite residual extension''.
\end{remarks}

\begin{lemma}
\label{III.1.3}
If $B$ is formally smooth over~$A$, then $B$ is flat over~$A$.
\end{lemma}

Since flatness is invariant under completion, we may suppose that~$A$ and~$B$ are complete. Since flatness is invariant under a local flat (hence faithfully flat) extension of the base ring, Definition~\ref{III.1.1} reduces us to the case where~$B$ is a formal power-series algebra over~$A$. But then, as an $A$-module, $B$ is isomorphic to a product of $A$-modules isomorphic to~$A$, hence (the base ring~$A$ being noetherian) is $A$-flat as a product of flat $A$-modules.

\medskip
Let us place ourselves under the conditions of~\ref{III.1.1}. Since the residual extensions of the local components of~$B'$ over~$A'$ are trivial, it follows that $L\otimes_k k'$ is an artinian $k'$-algebra whose local components have trivial residual extensions (where~$L$, $k$, $k'$ are the residue fields of~$A$, $B$, $A'$). This necessary condition for the finite free extension~$A'$ to satisfy the condition stated in~\ref{III.1.1} is also sufficient, as follows at once from~\ref{III.1.4}~(i) and~\ref{III.1.5} below.

\begin{proposition}
\label{III.1.4}
Let $A\to B$ be a local homomorphism of local rings, with finite residual extension, and let $A'$ be a finite local $A$-algebra over~$A$, so that $B'=B\otimes_A A'$ is finite over~$B$, hence a semi-local ring, also noetherian. \textup{(i)} If $B$ is formally smooth over~$A$, then the localizations of~$B$ at its maximal ideals are formally smooth over~$A'$. \textup{(ii)} The converse is true if~$A'$ is free over~$A$.
\end{proposition}

We are immediately reduced to the case where~$A$, $B$ are complete.

(i) Let~$A''$ be a finite free local extension of~$A$ such that the local components of~$B''=B\otimes_A A''$ are formal power-series algebras over~$A''$. Extending scalars $A''\to A''\otimes_A A'\to A'''$, where~$A'''$ is a local component of~$A''\otimes_A A'$, one sees that the local components of~$B''\otimes_{A''} A'''=B\otimes_A A'''$ are formal power-series algebras over~$A'''$. But one also has
\[
B\otimes_A A'''=(B\otimes_A A')\otimes_{A'} A'''=B'\otimes_{A'} A''',
\]
while, since~$A''$ is free over~$A$, $A''\otimes_A A'$ is free over~$A'$, and consequently so is~$A'''$, which is a direct factor of it. This proves that~$B'$ is formally smooth over~$A'$.

(ii) Let~$A''$ be a finite free local $A'$-algebra such that the local components of~$B'\otimes_{A'} A''=B\otimes_A A''$ are formal power-series algebras over~$A''$. Since~$A'$ is free over~$A$, so is~$A''$, and hence~$B$ is formally \emph{smooth} over~$A$.

\begin{proposition}
\label{III.1.5}
Let $A\to B$ be a local homomorphism of local rings, with \emph{trivial} residual extension. For~$B$ to be formally smooth over~$A$, it is necessary and sufficient that~$\overline{B}$ be isomorphic to a formal power-series algebra over~$\overline{A}$.
\end{proposition}

Only the necessity has to be proved, and we may suppose that~$A$ and~$B$ are complete. Let~$\goth{m}$ ($\goth{n}$) be the maximal ideal of~$A$ ($B$), and let~$t_1,\dots,t_n$ be elements of~$\goth{n}$ defining a basis of the vector space
\[
(\goth{n}/\goth{n}^2)/\operatorname{Im}(\goth{m}/\goth{m}^2)=\goth{n}/(\goth{n}^2+\goth{m}B).
\]
These elements therefore define a homomorphism of local $A$-algebras
\[
B_1=A[[t_1,\dots,t_n]]\to B.
\]
We prove that it is an isomorphism. It suffices to prove that for every power~$\goth{m}^q$ of~$\goth{m}$, one obtains an isomorphism after reducing modulo~$\goth{m}^q$ (since~$B_1$ and~$B$ are the projective limits of the corresponding rings reduced modulo~$\goth{m}^q$, with~$q$ variable). Since~$B$ and~$B_1$ are flat $A$-modules, the graded rings associated with the $\goth{m}$-adic filtration are obtained by tensoring over~$k=A/\goth{m}$ with $\gr(A)$ the rings $B_1/\goth{m}B_1$ respectively $B/\goth{m}B$. We are thus reduced to showing that $B_1/\goth{m}B_1\to B/\goth{m}B$ is an isomorphism. In view of~\ref{III.1.3}, we are thereby reduced to the case where~$A$ is a \emph{field}~$k$. On the other hand, if~$A'$ is a finite free local $A$-algebra such that~$B\otimes_A A'$ is a formal power-series algebra over~$A'$ (note that this algebra is local, since the residual extension of~$B$ over~$A$ is trivial), to prove that~$B_1\to B$ is an isomorphism, it suffices to prove that~$B_1\otimes_A A'\to B\otimes_A A'$ is one. This reduces us to the case where~$B$ is already a formal power-series algebra (this reduction should have been made first, before reducing to the case of a base field). But then~$B$ is a regular local ring with coefficient field~$k$, and it is well known (and immediate by considering the graded rings associated with the $\goth{n}_1$-adic and~$\goth{n}$-adic filtrations on~$B_1$ and~$B$) that~$B_1\to B$ is an isomorphism. This completes the proof.

\begin{corollary}
\label{III.1.6}
If~$B$ is formally smooth over~$A$, then there exists a finite local $A$-algebra~$A'$ such that the local components of
\[
\overline{B}\otimes_{\overline{A}}\overline{A'}=\overline{(B\otimes_A A')}
\]
are isomorphic to formal power-series algebras over~$\overline{A'}$.
\end{corollary}
Indeed, if~$L/k$ is the residual extension of~$B/A$, consider an extension~$k'/k$ such that the residual extensions in the $k'$-algebra~$L\otimes_k k'$ are trivial. Let~$A'$ be an algebra finite and free over~$A$ such that~$A'/\goth{m}A'=k'$ (one knows that such an algebra exists, for example by reducing step by step to the case where~$k'/k$ is monogenic, and then lifting to~$A$ the coefficients of the minimal polynomial of a generator of~$k'$ over~$k$). It is local. Then~$B\otimes_A A'$ has, at its maximal ideals, trivial residual extensions over that~$k'$ of~$A'$, and the conclusion follows with the help of~\ref{III.1.5}.

\begin{corollary}
\label{III.1.7}
Let~$A\to B$ be a local homomorphism of local rings. For~$B$ to be formally smooth over~$A$, it is necessary and sufficient that~$B$ be flat over~$A$ and that~$B/\goth{m}B$ be formally smooth over~$k=A/\goth{m}$.
\end{corollary}
Making a suitable finite free local extension~$A'$ of~$A$ and using~\ref{III.1.4}~(ii), we are reduced to the case where the residual extension of~$B/A$ is trivial. We know moreover by~\ref{III.1.4}~(i) and~\ref{III.1.3} that the stated conditions are necessary. For the sufficiency, it suffices to observe that the proof of~\ref{III.1.5} proves, under the hypotheses made here, that~$B$ is a formal power-series algebra over~$A$ (supposing~$A$ and~$B$ complete, which is permissible).

\begin{remark}
\label{III.1.8}
It would not be difficult to develop, for formally smooth homomorphisms, the analogue of all the properties of smooth morphisms studied in~\SourceRef{II}{II}. For the differential properties, however, this requires a modification of the usual definition of Kähler differentials (\cf I~\SourceRef{I.1}{1}), with completed tensor products replacing ordinary tensor products. We shall content ourselves with evoking these abysses here, what precedes being sufficient for our purpose.

It remains to make the link between formal smoothness and the notion of smoothness developed in~\SourceRef{II}{II} (and of which we have not yet made any use):
\end{remark}

\begin{proposition}
\label{III.1.9}
Let $A\to B$ be a local homomorphism, with~$B$ a localization of an $A$-algebra of finite type. For~$B$ to be smooth over~$A$, it is necessary and sufficient that it be formally smooth over~$A$.
\end{proposition}

Using~\ref{III.1.7} and~\SourceRef{II.2.1}{II.2.1}, we are reduced to the case where~$A$ is a field.

Using~\ref{III.1.4}~(ii) and~\SourceRef{II.4.13}{II.4.13}, a suitable extension~$k'$ of~$k$ reduces us to the case where the residual extension for~$B/k$ is trivial. In view of~\ref{III.1.5} (respectively~\SourceRef{II.5.2}{II.5.2}), $B$ is then smooth over~$k$ (respectively formally \emph{smooth over}~$k$) if and only if~$B$ is a regular local ring (respectively its completion is a formal power-series algebra over~$k$). But it is well known that these two conditions are equivalent (the residual extension being trivial).
